\documentclass[11pt,a4paper]{moderncv}

% Estilo Harvard
\moderncvstyle{classic} % puedes probar 'banking' o 'casual'
\moderncvcolor{blue}
\usepackage[scale=0.85]{geometry}

\name{Andrés}{López}
\title{Data Scientist \& Físico}
\address{Ciudad de México, México}{}{}
\phone[mobile]{+52~1~55~6543~4108}
\email{anderiotti@ciencias.unam.mx}
\social[github]{anderiotti}
\social[linkedin]{Andres Lopez}

\begin{document}
\makecvtitle

\section{Perfil}
Estudiante de Ciencia de Datos y Física con orientación a \textbf{Quantum Machine Learning}.
Apasionado por la investigación computacional, el análisis estadístico con alto formalismo matemático y computacional, y el desarrollo de modelos predictivos de alto impacto. 
Capacidad demostrada para integrar teoría cuántica, algoritmos de aprendizaje automático y sistemas de información. 
Enfocado en transformar conocimiento científico en soluciones analíticas y modelos predictivos que optimicen procesos y decisiones basadas en datos.


\section{Educación}
\cventry{2021--Presente}{Licenciatura en Física}{Universidad Nacional Autónoma de México}{Ciudad de México}{}{
Fuerte enfoque en modelado de sistemas complejos, mecánica cuántica, análisis numérico y procesamiento de señales. 
Experiencia en análisis de datos y machine learning en laboratorios de física experimental. 
\textbf{Promedio: 9.5/10.}}

\cventry{2021--Presente}{Licenciatura en Ciencia de Datos}{Universidad Nacional Autónoma de México}{Ciudad de México}{}{
Especialización en optimización, teoría de grafos, estadística multivariada y análisis funcional. 
Desarrollo de modelos de regresión, clasificación, clustering y análisis de series de tiempo, así como teoría general en bases de datos estructuradas y no estructuradas. 
\textbf{Promedio: 9.6/10.}}


\section{Experiencia Relevante}
\cventry{2024--2025}{Analista de Ciencia de Datos}{Grupo de Investigación Estudiantil en Facultad de Ciencias, UNAM}{Ciudad de México}{}{
Desarrollo de modelos predictivos en \textbf{Python} y \textbf{R} para series temporales de experimentos físicos.  
Limpieza de datos con \emph{SQL} y planificación de sistemas OLTP en \emph{Power BI}.}
\cventry{2025}{(IBM) Qiskit Global Summer School 2025 on Quantum Computing}{Quantum Badge}{}{}{
Exploración formal en la intersección de la computación clásica y cuántica, aplicando técnicas de machine learning para la mitigación de errores por ruido o decoherencia cuántica, desarrollo de modelos de Quantum Machine Learning (QML) para la clasificación de datos y programación en computadoras cuánticas reales.}
\cventry{2024}{Proyecto de Análisis Magnetohidrodinámico en Fluidos Reoscópicos}{Laboratorio de Hidrodinámica y Turbulencia, UNAM}{Ciudad de México}{}{Lider en estudio para modelar cómo la topología del campo electromagnético induce la transición de un flujo predecible a un régimen de dinámica caótica, cuantificando la sensibilidad del sistema a las condiciones geométricas iniciales.}

\section{Competencias Técnicas}
\cvitem{Lenguajes}{Python, Julia, R, C, C++, SQL}
\cvitem{Machine Learning}{scikit-learn, TensorFlow (Keras), PyTorch, modelos de clasificación y regresión, reducción de dimensionalidad, aprendizaje no supervisado, series temporales}
\cvitem{Quantum Computing}{Qiskit, PennyLane, fundamentos de computación cuántica, modelos híbridos QML}
\cvitem{Data Analysis}{Pandas, NumPy, Matplotlib, Seaborn, análisis exploratorio, optimización numérica}
\cvitem{Business Intelligence}{Power BI, Tableau, Excel avanzado (PowerQuery, tablas dinámicas), Power Designer (modelado de datos)}
\cvitem{Bases de Datos}{Diseño y consultas en SQL, modelado relacional, SQLAlchemy}
\cvitem{Nube y Dev Tools}{Git, Linux, Docker básico, AWS, VirtualBox}


\section{Idiomas}
\cvitemwithcomment{Inglés}{Avanzado}{Lectura y conversación fluida}
\cvitemwithcomment{Italiano}{Elemental}{Lectura aceptable}

\section{Habilidades Blandas}
\cvitem{Pensamiento analítico}{Capacidad destacada para descomponer problemas complejos en estructuras lógicas y formular soluciones precisas, integrando perspectivas matemáticas, físicas y computacionales.}
\cvitem{Rigor conceptual}{Tendencia natural a mantener un alto nivel de precisión teórica en el diseño e interpretación de modelos, garantizando coherencia entre los fundamentos físicos y la implementación algorítmica.}
\cvitem{Comunicación técnica}{Habilidad para traducir conceptos abstractos de alta complejidad en explicaciones claras y estructuradas, adaptadas a audiencias multidisciplinarias.}
\cvitem{Colaboración interdisciplinaria}{Experiencia trabajando en equipos mixtos de físicos, matemáticos, matemáticos aplicados, actuarios, ingenieros y analistas de datos, promoviendo la integración metodológica entre el razonamiento teórico y la aplicación práctica.}
\cvitem{Aprendizaje autónomo}{Alta capacidad de autogestión en la adquisición de nuevos lenguajes de programación, bibliotecas científicas y paradigmas teóricos, así como del formalismo matemático puro, priorizando la comprensión profunda sobre la memorización superficial.}
\cvitem{Gestión del tiempo}{Eficiencia en la planificación de tareas simultáneas de investigación, programación y análisis, bajo entornos de alta demanda cognitiva.}
\cvitem{Disciplina y constancia}{Miembro activo del equipo representativo de powerlifting de la UNAM, combinando entrenamiento competitivo de alto rendimiento con exigencias académicas avanzadas, demostrando resiliencia, compromiso y control mental bajo presión.}
\cvitem{Excelencia académica}{Estudiante destacado en Física y Ciencia de Datos, con desempeño sobresaliente en ambas carreras, evidenciando una ética de trabajo rigurosa, capacidad de concentración prolongada y una orientación constante hacia la mejora continua.}



\clearpage
\end{document}
